\documentclass[11pt]{article}
\usepackage{amsmath,amssymb,amsthm,hyperref}
\newtheorem{proposition}{Proposition}
\title{Vector Cauchy Lookup and Admissible Two-Head Products}
\author{Research checkpoint, branch ar/027ac8030b43}
\date{November 24, 2026}
\begin{document}
\maketitle

\paragraph{Status.}
The two-head product construction and simultaneous preservation of strict
signs are kernel checked in \texttt{Results/LargeDomainProducts.lean}.
The lookup count and budget analysis below are written arguments using
Warren's theorem, not Lean declarations. There is no improved unconditional
bound on the original worst-case head complexity in this note. The model,
admissibility conditions, and definition of that complexity are unchanged.

\section{A concrete family-dependent lookup interface}

Fix distinct tokens $t_1,\ldots,t_q$ and exterior poles
$D_1,\ldots,D_q$, with $D_u>t_z$. Given an arbitrary table
$T:[r]\times[s]\to[q]$, seek a single score on the three one-hot axes
\[
 S_T(u,v,j)=c+\sum_{h=1}^H
 \frac{a_{h0}+a^U_{hu}+a^V_{hv}+a^J_{hj}}
      {b_{h0}+b^U_{hu}+b^V_{hv}+b^J_{hj}}
 \approx\frac1{D_u-t_{T(v,j)}}.
 \tag{1}
\]
Every displayed coefficient may depend jointly on the complete table $T$.
The same heads and coefficients must then work simultaneously for all
$u,v,j$. Each denominator is an affine function on the full Boolean cube
of dimension $q+r+s$ and must be positive everywhere on that cube, with
the required common slope orientation. Approximation means arbitrarily
small uniform error, because the outer finite-plane decoder may have
arbitrarily ill-conditioned fixed residues. A proof of signs alone is not
a proof of this lookup interface.

The naive expansion into $q$ token-indicator matrices uses
$O(qr/\log r)$ heads if each indicator matrix is compiled separately.
A possible efficient target is $O(r\log q/\log r)$. This interface is more
general than a fixed row dictionary: the row slopes and numerators are
allowed to depend on all independently chosen column tables together.
A separation-rank lower bound does not apply to an arbitrary sum in (1),
whose denominators couple the three axes.

\section{A nonlinear product that avoids a numerator-rank charge}

Let $A,D,f$ be affine functions on a Boolean cube, with $D>0$ everywhere
and every slope of $D$ nonpositive. Write $f(x)=f_0+\sum_i f_i x_i$ and put
\[
 G(x)=1+|f_0|+\sum_i(1+2|f_i|)-\sum_i(1+|f_i|)x_i.
\]
Both $G$ and $G+f$ are strictly positive on the whole cube; their slopes
are strictly negative. For every $t>0$, the two heads
\[
 Q_t(x)=\frac{tA(x)}{D(x)+G(x)/t}
       -\frac{tA(x)}{D(x)+(G(x)+f(x))/t}
\]
therefore have positive denominators everywhere and strictly negative
slopes. Their numerators remain affine. Direct cancellation gives
\[
 Q_t(x)=\frac{f(x)A(x)}
 {(D(x)+G(x)/t)(D(x)+(G(x)+f(x))/t)}
 \longrightarrow\frac{f(x)A(x)}{D(x)^2}.
 \tag{2}
\]
The convergence is simultaneous on every finite set. Consequently, if a
prior score $P$ has the desired strict signed margins after adding
$fA/D^2$, one sufficiently large common $t$ preserves all those margins
with the two admissible heads instead. Every zero of $A$ is an exact zero
of both new heads for every $t$.

The checked declarations in namespace \texttt{HeadComplexity.LargeDomain}
are \texttt{exists\_strict\_reverse\_affine\_product\_pair},
\texttt{strict\_reverse\_affine\_product\_pair\_limit}, and
\texttt{affine\_product\_pair\_preserves\_signs}. The two supporting guard
positivity declarations are also checked. The construction uses no
composition layer, and its two heads satisfy positivity on multihot inputs
as well as on the local one-hot domain.

This is a potential tool for (1), not a solution. It multiplies an affine
selector by a squared reciprocal kernel. A full inner compiler based on
squared Cauchy kernels, its common parameter choices, and its complete
axis repair still need separate proofs. Merely applying (2) to an ordinary
reciprocal decoder changes its kernel and cannot be treated as a black-box
multiplication theorem for that decoder.

\section{Why an efficient three-axis lookup still does not multiply savings}

First, the $\log q$ cost in the optimistic target is necessary in the
polynomial regime. Restrict (1) to one fixed outer pole $u_0$. The $q$
values $1/(D_{u_0}-t_z)$ are distinct, so $q-1$ intermediate thresholds
recover $T(v,j)$ from any sufficiently accurate approximation. Thus all
$q^{rs}$ color tables give distinct strict sign vectors at
$m=(q-1)rs$ threshold tests.

For $H$ heads, use $p=1+2H(r+s+1)$ real parameters after restriction to
$u_0$. Clearing the $H$ positive denominators makes each threshold test a
polynomial of degree at most $H+1$. Provided $m\geq p$, Warren's theorem
therefore gives the necessary inequality
\[
 rs\log_2q\leq p\log_2\left(\frac{4e(H+1)(q-1)rs}{p}\right).
 \tag{3}
\]
The alternative $p>m$ already gives a stronger lower bound in the regimes
at issue. If $s$ is comparable to $r$, $2\leq q\leq r^B$, and
$H\leq r^C$ for fixed constants, (3) implies
\[
 H=\Omega_{B,C}\left(\frac{r\log q}{\log r}\right).
 \tag{4}
\]
This proof allows label-dependent row coefficients and arbitrary
interference among the heads. It excludes a universal lookup bound
$o(r\log q/\log r)$ in this regime, not a general Boolean sign decoder
which avoids recovering the lookup values.

Now take an outer finite plane with $P=q^2$ memory points and
$a=\log_2q$ bit queries per point. Across inner axes of sizes $r,s$, the
full tensor has $M=P a r s$ independent Boolean labels. Even an optimal
lookup of cost $O(r\log q/\log r)$ at every outer memory point gives
\[
 H=O\left(\frac{P r\log q}{\log r}\right)
   =O\left(\frac{P a r}{\log r}\right),\qquad
 \frac{M}{H}=\Theta(s\log r).
 \tag{5}
\]
The outer $a$-bit saving is consumed by the $\log q$ cost of the inner
lookup. For $P a,r,s$ of order at most the ambient dimension, this is the
inherited $n\log n$ inputs-per-head scale. This is consistent with the
independent-cell size test: a three-axis domain is polynomial in $n$ and
cannot support any unbounded improvement over that scale.

A growing number of axes remains possible. For $d$ axes of size about
$n/d$, the input count is $(n/d)^d$ and the logarithm is
$d\log(n/d)$. A successful joint decoder must exploit this whole
logarithm while keeping the cover density and boundary repairs controlled.
Applying independent vector lookups recursively only supplies the bottom
axis logarithm; no multiplicative logarithmic gain follows from (1).

\paragraph{Analytic source.}
The sign-count input is Warren's strict polynomial sign-pattern bound;
a convenient primary statement is Theorem 2.1 of
\href{https://www.cs.ox.ac.uk/people/paul.goldberg/papers/goldbergjerrum.pdf}
{Goldberg and Jerrum, Bounding the Vapnik--Chervonenkis Dimension of
Concept Classes Parameterized by Real Numbers}.
No analytic sign-count theorem is introduced as a Lean axiom.
\end{document}
